\begin{table}[H]
\centering
\caption*{\textbf{Quadro 4 - Exemplos de montagem do menu de tools}}
\footnotesize
\renewcommand{\arraystretch}{1.12}
\begin{tabularx}{\linewidth}{|>{\raggedright\arraybackslash}p{0.27\linewidth}|>{\raggedright\arraybackslash}p{0.36\linewidth}|Y|}
\hline
\textbf{Configuração} & \textbf{Bundle} & \textbf{Menu resultante} \\
\hline
$T_{\mathrm{base}} = \{\texttt{calculator}\}$ &
vazio & $\{\texttt{calculator}\}$ \\
\hline
$T_{\mathrm{base}} = \{\texttt{calculator}\}$ &
\texttt{latest-document-reading} declara \texttt{list} e \texttt{read} &
$\{\texttt{calculator},\texttt{list},\texttt{read}\}$ \\
\hline
$T_{\mathrm{base}} = \{\texttt{calculator}\}$ &
\texttt{slack-message} declara \texttt{slack\_list\_channels} e
\texttt{slack\_send\_message} & união das três tools \\
\hline
YOLO máximo: $T_{\mathrm{base}} = T_{\mathrm{registry}}$ &
qualquer & todas as tools registradas \\
\hline
\end{tabularx}
\par\smallskip\raggedright\small Fonte: elaboração própria (2026).
\end{table}
